\section{Introduction}\label{sec:intro}

At fixed information, a stronger incumbent reduces a prospective challenger's expected return from entry into a
takeover contest. The challenger must pay preparation costs before it can submit an executable offer, and a
stronger rival lowers what that preparation earns \citep{Fishman1988,HirshleiferPng1989}. I study a force in
the opposite direction. Competition can make the target's stock more informative about the challenger. When a
potential buyer observes that price before it commits to acquisition preparation, the information effect can
outweigh the loss of acquisition profits. A stronger incumbent then attracts another buyer and raises the
probability that a challenger from the high-value class acquires the target.

Consider a listed target with a public sale opportunity. One buyer is already prepared, and another has
not committed to the diligence that an executable proposal requires. Section~\ref{sec:interval} describes the
decision interval in which that buyer decides while the target's stock trades. Incumbent strength is the public
distribution of the prepared bidder's value, and a stronger incumbent shifts it upward. No offer is announced
before the challenger decides.

The mechanism runs through two claims on the same acquisition surplus. The challenger holds the acquisition
claim: what it keeps if it prepares, bids, and wins. Target shareholders hold the traded claim: what the cash
second-price auction pays them. The challenger's value comes from a high-value or a low-value quality class, and the
classes fix the expected value gap between them. A stronger incumbent outbids the challenger more often and then pays
the challenger's bid, so more of the gap reaches target proceeds. A challenger that wins pays more, so less of the gap
remains as acquisition profit at every belief.

Target shares therefore become more sensitive to information about the challenger, while the challenger's own return
to preparation falls. An outside specialist investor who knows the challenger's class trades on that sensitivity. It
holds no net position in the target. Target shareholders gain when the price recruits entry, because a holder through
the sale receives sale proceeds, not the interim price.

The trading stage is one order slot, with the alternative-arrival structure of \citet{DowGorton1997}. With some
probability the informed investor fills it. Otherwise a liquidity trader fills it with a sale, a purchase, or no
order, and no one observes who placed the order. A liquidity order in the same direction is always possible, so no
order reveals the class. The posterior after every order stays inside fixed bounds under every investor strategy, and
the investor keeps a residual advantage after every order. A larger gap in target proceeds can make that advantage
worth the cost of trade.

The main result compares a weak-incumbent economy with a strong-incumbent economy on a nonempty open set of
primitives. In each economy the trading outcome and the on-path preparation outcome are unique. With the weak
incumbent, the investor's residual advantage cannot cover the trading cost, the price reveals nothing, and only
a zero-cost challenger prepares. With the strong incumbent, the investor buys after good news and sells after
bad news, and a purchase brings in a high-cost challenger. Entry rises even though acquisition profit falls at
every fixed belief. The probability that a challenger from the high-value class owns the target rises with it.

Market makers anticipate the preparation that each price induces, so the reversal holds under rational
pricing. Uniqueness covers every mixed order profile on the menu \(\{-1,0,1\}\) and every unilateral deviation.
The posterior bounds hold for every profile, and they deliver this uniqueness. The change in information does
the work. If the investor's orders stay fixed at their informative level while the incumbent strengthens, entry
does not rise (Section~\ref{sec:benchmark}).

In the benchmark of Section~\ref{sec:benchmark}, entry rises from \reg{entry_r0} at the weak incumbent to
\reg{entry_r1_B} at the strong incumbent. High-class ownership rises from \reg{own_r0} to \reg{own_r1_B}.
Over the same change, gross acquisition profit at the prior falls from \reg{Bprior_r0} to
\reg{Bprior_r1_B}, and the target-payoff spread rises from \reg{Delta_r0} to \reg{Delta_r1}. Against a
strong enough incumbent, even the most favorable price cannot justify expensive preparation. Entry then returns
to its floor while the price stays informative, so entry is non-monotone in strength. In the benchmark,
high-cost preparation stops above the strength \(r_C\approx\reg{rC_B}\).

The hypotheses of the main result are strict inequalities, so they define an open set of primitives. At the
benchmark auction primitives, they hold exactly when \(\reg{u_lo_frac}<u<\reg{u_hi_frac}\). Here \(u\) is the informed
presence mass relative to the liquidity mass in one direction. The number \(1+u\) is the largest likelihood ratio
that one order can carry. Equivalently, the top posterior at the prior \(1/2\) lies between the preparation thresholds of the
strong and the very strong incumbent, \(\reg{M_lo_frac}\) and \(\reg{M_hi_frac}\). Section~\ref{sec:benchmark} reports
this region and the range of priors as a table.

Section~\ref{sec:coexistence} gives the complete equilibrium set between and beyond the benchmark strengths.
Informative and uninformative equilibria coexist on an interval of strengths \([r_A,r_N]\), with
\(r_A\approx\reg{rA}\) and \(r_N\approx\reg{rN}\). On the informative
branch, the high-class investor first mixes between a purchase and abstention, and entry rises strictly with
strength. From the strength \(r_F\approx\reg{rF_B}\), the investor buys for sure, and entry stays flat on the rest
of the interval. Entry therefore rises weakly along the informative branch and strictly where the investor mixes. The
same start of coexistence holds for a window of priors around the benchmark.

Section~\ref{sec:research} lets the investor pay for research that makes it informed. The investor chooses the
probability of success up to a cap, at a quadratic cost. The unique outcome has no research with the weak incumbent
and research at the cap with the strong incumbent. At intermediate strengths, research and entry rise together. At the
strengths \reg{str_r155}, \reg{str_r160} and \reg{str_r165}, research rises from \reg{res_lambda_r155} through
\reg{res_lambda_r160} to \reg{res_lambda_r165}, and entry from \reg{res_entry_r155} through \reg{res_entry_r160} to
\reg{res_entry_r165}. Along this branch the comovement holds for every liquidity probability of the model, and an
equilibrium without research exists beside it.

The reversal holds at some incumbent strengths and costs for every prior and every pair of classes that differ
above the reserve. It needs the high class to have stochastically higher values and both class supports to lie
above the reserve (Section~\ref{sec:classes}). With more than two classes, it holds under a separation
condition. The rise of entry survives other laws of preparation costs, with or without atoms
(Section~\ref{sec:costs}). It also survives private signals that the buyer and the investor hold about different
facts (Section~\ref{sec:signals}). In an example with its own costs and strengths, the buyer's signal is more
accurate than the investor's, and the rise needs both signals.

In the strong-incumbent economy, access to prices raises expected target proceeds and acquisition surplus net
of preparation costs (Section~\ref{sec:access}). Sale terms shape both what a buyer pays when it wins and what the
stock price can reveal before anyone prepares. A reserve that excludes the low-value class makes the price
informative against both incumbents, and entry is then flat in incumbent strength (Section~\ref{sec:reserves}).
An alternative bargaining institution shows how the payment rule matters (Section~\ref{sec:payment}). The effect
of competition on the target-payoff spread depends on the weight the rule places on the runner-up's value.

The model also gives tests (Section~\ref{sec:empirical}). The purchase impact exceeds the sale impact exactly when a
purchase recruits the buyer. In the benchmark, this entry component is \reg{ba_share_r1_B} of the spread between
purchase and sale prices. Against a strong prepared bidder, orders no longer predict entry, but they still predict
which bidder wins. Given the entry floor and the target-payoff spread, the sale impact and the frequency of orders
identify the informed share.

The paper belongs to the literature on the real effects of financial markets \citep{BondEdmansGoldstein2012}.
The closest antecedent is \citet{DowGoldsteinGuembel2017}, where a firm's investment decision feeds back into
the incentive to produce information about its stock. Here the sale rule divides acquisition surplus between a
traded claim and a buyer that decides whether to enter. \citet{EdmansGoldsteinJiang2015} show that corrective
real decisions can discourage trade on bad news under rational pricing. Here the real decision is the buyer's
entry itself.

Firms learn from peers' stock prices \citep{FoucaultFresard2014}, and target prices affect takeover activity
\citep{EdmansGoldsteinJiang2012}. In this model, a prospective challenger learns about its own acquisition value
from the target's price before it commits. Bidders also learn from announcement returns when they
decide whether to complete a deal \citep{Luo2005}. That decision comes after this entry margin.

The auction-entry literature makes the bidder pool endogenous
\citep{Samuelson1985,McAfeeMcMillan1987,LevinSmith1994,GentryStroup2019}. It also shows how selective entry
affects the choice of sale procedure \citep{RobertsSweeting2013}. The challenger here faces selective entry in
the terms of that literature. Its pre-entry information indexes the distribution of its value, and the payment
of the entry cost reveals the value \citep{RobertsSweeting2013}. One limit of this framework is a bidder that
knows only the common value distribution before entry \citep{LevinSmith1994}. The other limit is a bidder that
already knows its value \citep{Samuelson1985}.

Here the pre-entry information is the target's stock price. That price is public, and it moves with the orders
of an investor who observes the challenger's quality class. I keep the direct deterrence force of the
auction-entry literature and add a market that operates before preparation. Auction formats also affect
bidders' incentives to acquire information \citep{Persico2000}. The informed trader here is outside the
auction and trades a payoff that differs from the entrant's profit.

Recent work studies how bidders learn about their own values and about competitors \citep{PernoudGleyze2026}. It also studies
post-auction feedback in security-payment design \citep{LiuBernhardt2022} and bidder-pool choice with correlated
values \citep{CarlinEtAl2026}. In takeovers, \citet{BettonEtAl2014} study negotiations with stock-market
feedback, and \citet{LinMaYangZhu2025} model payment choice and trade within an initiated deal. I hold cash
consideration fixed and study entry.

Arbitrageurs' positions can affect tendering \citep{CornelliLi2002}. The investor here has information about a
prospective acquirer and no role in tendering. The model abstracts from toeholds and dispersed-shareholder free
riding \citep{BulowHuangKlemperer1999,GrossmanHart1980}. Section~\ref{sec:research} makes the investor's
information the result of costly research, with the alternative-arrival structure of \citet{DowGorton1997}.
\citet{GoldsteinGuembel2008} study a related channel, manipulation through the real decision that responds to
the price.

An informed trader can also hold a stake before a control decision. Liquidity trade then lets an acquirer
conceal share purchases before a bid \citep{KyleVila1991}. Large holders trade on their information
\citep{Edmans2009}, monitor, or intervene as prices feed back to firm decisions
\citetext{\citealp[Section~II.E]{EdmansGoldsteinJiang2015}; \citealp{Maug1998}; \citealp{BackEtAl2018}}. The
investor in this paper holds no stake, seeks no control, and earns only the trading residual.

Section~\ref{sec:model} describes the decision interval and sets out the model. Section~\ref{sec:payoffs}
derives the two returns to information and what the challenger learns from the price.
Section~\ref{sec:results} states the main result, the benchmark with its fixed-information control, and the
coexistence of equilibria. Section~\ref{sec:research} adds costly research. Section~\ref{sec:robustness} covers
robustness, Section~\ref{sec:welfare} welfare and sale terms, and Section~\ref{sec:empirical} empirical
implications, and Section~\ref{sec:conclusion} concludes. Appendix~\ref{app:proofs} proves every result and states the
results of Sections~\ref{sec:robustness} to~\ref{sec:empirical} in full.
